% !TEX root = ../main.tex
\lecture{6}{2026-10-24}{Selflearning}

\subsection{Einkommens und Substitutionseffekt}
Bisher konnte anhand einer Nutzenfunktion, Warenpreisen im Bündel und der Budget Bedingung, die Optimale Verteilung (optimale Menge beider Güter, die der Konsument nachfragen sollte) hergeleitet werden.
Es stellt sich die frage des Zusammenhangs und Effekte von Änderungen des Budgets und Preisen, die vorher als konstanten betrachtet wurden, auf die Optimale Verteilung (Marschall'sche Nachfrage).

\subsubsection{Einkommens-Konsum-Kurve}
Mit steigendem Budget und keine Veränderungen der Preise, können Akteure von allen Güter eines Bündels mehr konsumieren, jedoch in noch verhältnismässig gleicher Mengen. Die Nichtsättigung führt dazu, dass der Nutzen steigt, wenn die Möglichkeit besteht Bündel grösserer Beträge zu wählen.

Es entstehen höhere oder niedere Indifferenzkurven auf der Nutzenfunktion, die wiederum die Bedingung die Optimalitätsbedingung bei neuem Budget erfüllen. Die Optimumpunkte der Indifferenzkurven mit den entsprechenden Budgetbedingungen liegen auf einer Gerade (Einkommenseffekt). Diese Gerade wird als \textbf{Einkommens-Konsum} oder \textbf{Engel-Konsum-Gerade} bezeichnet. Diese Funktion zeigt wie das Budget steigen muss mit steigendem Optimaler menge an einer Ware des Bündels, oder umgekehrt zeigt welche Optimale Menge der Ware dem steigendem Budget entspricht.


\subsubsection{Preis-Konsum-Kurve \& Nachfrage}
Budget bleibt konstant während Preise der Waren sich ändern. Gleichzeitige Veränderungen der Preise aller Waren in gleichem Ausmass führen zum gleichen Effekt als wäre das Budget verändert und Preise gleichgeblieben. Kombinationen würden verhältnismässig gleich bleiben mit mehr oder weniger von beiden Waren.
Bei der einseitigen oder unterschiedlichen Veränderung des Preises der Waren wird die Budgetgerade mit den Güterkombinationen nach innen oder aussen rotiert. Sinkt der Preis eines Gutes W (x-achse), kann mehr von ihm konsumiert werden bei gleichem Budget Schnittpunkt bewegt sich nach rechts.
Die Verbindung der Optimumspunkte bei Preisveränderungen nennt sich Preis-Konsum-Kurve und zeigt welche Optimale Menge bei Veränderungen im Preis konsumiert werden können.
Resultiert aus dem einfügen der Bedingung der Preis Veränderung in der Optimalitätsbedingung.
Die resultierende Gerade stellt für ein Individuum bei Optimale Konsummenge einem entsprechenden Preis (Zahlungsbereitschaft), diese entspricht die\textbf{Nachfrage-Kurve}

\subsubsection{Preisänderung: normale Güter}

\begin{definition}
    [normales Gut] ist eine Ware, dessen nachgefragten Menge mit steigender Kaufkraft der individuellen Akteure auch steigt. (Kinobesuche, Reisen, usw.)  
\end{definition}

\begin{itemize}
    \item Substitutionseffekt (SE)
    \begin{itemize}
        \item Wenn bei steigendem Preis eines Gutes (W), wird es durch eine höhere Menge des anderen günstigeren Gutes (F) im Bündel ersetzt. Negativ Effekt auf Nachfrage von W
        \item führt zu Bewegung des Optimumpunkts auf der Indifferenzgerade (dieser liegt jetzt unter oder über die originale Budgetgerade)
        \item \textbf{Immer Negativ}, das Steigen des Preises führt zu reduzierten Nachfrage des Produkts und das Sinken des Preises führt zu deren Erhöhung
    \end{itemize}
    \item Einkommenseffekt (EE)
    \begin{itemize}
        \item Durch Steigung des Gesamtpreises wird die Kaufkraft reduziert (kann sich weniger Produkte leisten für gleiche Geldmenge)
        \item führt zu Einkommenseffekt und parallele Verschiebung zu einer Kurve niedrigerem Nutzens
        \item so wird der durch SE verursachte neue Punkt wieder zum Optimum einer anderen Budget gerade
    \end{itemize}
    \item Gesamteffekt
    \begin{itemize}
        \item Setzt sich aus den beiden Effekten zusammen
        \item Ist bei normalen Gütern \textbf{immer Negativ}, weil Definition besagt, dass es Güter sind, die proportional zum Einkommen nachgefragt werden.
    \end{itemize}
\end{itemize}

\subsubsection{Preisänderung: inferiore Güter}

\begin{definition}
    [inferiores Gut] ist eine Ware, dessen nachgefragten Menge mit steigender Kaufkraft der individuellen Akteure sinkt/durch andere qualitativ hochwertigere Produkte ersetzt wird. (Fast-foot->Restaurants usw.)  
\end{definition}

\begin{itemize}
    \item Gesamteffekt
    \begin{itemize}
        \item SE und EE können gegensätzliche Wirkungen haben, weil der Konsum inferiorer Güter vom Einkommenseffekt positiv wirkt, während SE negativ. 
    \end{itemize}
\end{itemize}

\begin{definition}
    [Giffengüter] bei stark inferioren Güter steigt oder bleibt konstant die Nachfrage auch mit steigenden Preisen. (z.B für das Überleben notwendige Produkte, Brot, Wasser, etc.)
\end{definition}

\subsection{Perfekte Komplemente}

\begin{definition}
    [Perfekte Komplemente] sind zwei oder mehrere Güter, die nur Nutzen stiften, wenn beide im gleichen oder sonst fixen Verhältnis vorhanden sind. (z.B. Schuhe und Socken)
\end{definition}

\subsubsection{Eigenschaften}
\begin{itemize}
\item Nutzensfunktion ist pyramidenförmig
\end{itemize}

\subsubsection{Effekte auf Indifferenzkurven}

Sehen aus wie senkrechte "L" förmige geraden, weil aufgrund des festen Verhältnisses perfekter Komplemente keine mögliche Substitution gibt. Man könnte nicht weniger Socken haben als Schuhe, und mehr Schuhen wären nutzlos ohne entsprechende Menge an Socken.


\textbf{Nur Einkommenseffekt möglich}, weil dieser keinen Einfluss auf das Verhältnis der konsumierten Waren hat. Entweder mehr oder weniger von beiden Güter

\subsection{Perfekte Substitute}

\begin{definition}
    [Perfekte Substitute] sind Güter die sich in einem fixen Verhältnis ohne Veränderung des Nutzens austauschen lassen. Es herrscht Indifferenz. 
\end{definition}

\subsubsection{Eigenschaften}
\begin{itemize}
    \item Nicht strikt Konvex (unausgewogene extreme sind Möglich)
\end{itemize}

\subsubsection{Effekte Preiserhöhungen}
\begin{itemize}
    \item GE = SE
    \begin{itemize}
        \item Wenn Gut dessen Preis erhöht wurde nicht mehr konsumiert wird und dafür nur das andere
    \end{itemize}
    \item GE = EE
    \begin{itemize}
        \item Wenn trotz Preisveränderung keine Konsumveränderung gibt
        \item Nach wie vor keine Menge an Gut 
    \end{itemize}
\end{itemize}

\subsection{Elastizitäten}

\subsubsection{Motivation}
\begin{itemize}
    \item Wie müssen Preise gesetzt werden um Umsätze zu erhöhen anhand der Arten von Nachfrage?
    \item Wenn Preissenkung zu mehr Umsatz durch erhöhte Nachfrage führt ist diese Nachfrage elastisch
    \item Wenn Preiserhöhung zu mehr Umsatz führt ist Nachfrage unelastisch (z.B Wasser, Nahrung)
    \item Frage lässe sich nicht apriori beantworten, man müsse zum Anfang eine Option wählen und beobachten
\end{itemize}

\subsubsection{Berechnung}

Die Elastizität ist ein einheitloses Mass (Verhältnis), das durch das Ableiten der Preis-Nachfrage-Kurve in Prozentuale Änderungen in Y und in X relativ zu sich selbst berechnet wird.

\begin{equation}
    \eta_{Y,X}(X) = \frac{dY / Y}{dX / X} = \frac{f'(x)\cdot x}{f(x)} = f'(x) \cdot \frac{x}{y}
\end{equation}

dadurch lassen sich verschieden skalierte Preis-Nachfrage-Kurve sich miteinander vergleichen. Diese unterschiedliche Skalierung ist auf die unterschiedlichen Masseneinheiten der Waren zurückzuführen.

\subsubsection{Einkommenselastizität}
Wie ändert sich die Nachfrage, bei Veränderung des Einkommens/budgets und konstante Preise? Die Antwort dient zur Klassifizierung von Güter in:
\textbf{Kategorien;}
\begin{itemize}
    \item inferiore Güter ($ \eta < 0$)
    \item normale ($ 0 < \eta$)
    \begin{itemize}
        \item notwendige ($0 < \eta < 1$)
        \item Luxusgut ($ 1 < \eta $)
    \end{itemize}
\end{itemize}

\subsubsection{Ausgabenanteil}

Wie ändert sich die Ausgabe für eine Bestimmte ware in Relation zum Einkommen, wenn sich das Einkommen variiert?
\begin{equation}
    A(M) = \frac{P \cdot Q(M)}{M}
\end{equation}

\begin{equation}
    A'(M) = \frac{M\cdot P \cdot Q'(m) - P \cdot Q(M)}{M^2}
\end{equation}

\begin{itemize}
    \item Wenn Steigung $\le 0$ => normales Gut
    \item > 0 => Luxusgut
\end{itemize}

\subsubsection{(Eigen-) Preiselastizität}
Wie verändert sich die Nachfrage eines Gutes im Verhältnis zu seinem Preis? 

\begin{equation}
    \eta_{Q,P}(P) = \frac{dQ / Q}{dP / P} = Q'(P) \cdot \frac{P}{Q(P)}
\end{equation}
Wichtig für die Frage der Preissetzung wenn Preis die nachgefragte Menge beeinflusst und Umsatz von Menge und Preis pro Menge abhängig ist.
\textbf{TODO: Review Bedeutung von postive, negative Elastizität}

\subsubsection{Kreuzpreiselastizität (KPE)}
Misst wie die Änderungen im Preis eines Gutes die Nachfrage eines anderen beeinflussen (Substitute, Komplemente, oder keinen Einfluss)

\begin{itemize}
    \item Bei \textbf{Substitute} steigt die Nachfrage im Preis des anderen -> positive KP-Elastizität
    \item Bei \textbf{Komplemente} sinkt die Nachfrage im Preis des anderen -> negative KP-Elastizität
\end{itemize}